%! Parametric Functions — Unit 8, Lesson 8.1 — Guided Notes
\documentclass[10pt]{article}
\usepackage{precalculus-article}
\usepackage{precalculus-boxes}
\newcommand{\TallMath}[1]{$\displaystyle #1\rule[-1.4em]{0pt}{3.2em}$}

\begin{document}

\pageheader{Unit 8, Lesson 8.1}{Guided Notes}
\namedateperiod

\vspace{0.10in}

%----------------------------------------------------------------------
% Objectives
%----------------------------------------------------------------------
\begin{objectivebox}
\textbf{I will be able to\ldots}
\begin{enumerate}[leftmargin=1.5em, itemsep=4pt]
  \item \textbf{LO 4.1.A} --- Construct a table of values and a graph
    for a parametric function $f(t) = (x(t),\, y(t))$ by evaluating both
    component functions at selected values of the parameter $t$, plotting
    the resulting ordered pairs in order of increasing $t$, and
    interpreting domain restrictions as start and end points. \writeline
\end{enumerate}
\end{objectivebox}

\vspace{0.08in}

%----------------------------------------------------------------------
% Vocabulary
%----------------------------------------------------------------------
\begin{vocabbox}
\begin{description}[leftmargin=0pt, itemsep=6pt]
  \item[\termblanklong{parameter}]
    The \underline{\hspace{3cm}} variable $t$ in a parametric function;
    it controls both $x$ and $y$ at the same time.

  \item[\termblanklong{parametric function}]
    A function $f(t) = (x(t),\, y(t))$ in $\mathbb{R}^2$ consisting of
    \underline{\hspace{2cm}} component equations, both written in terms
    of a single variable $t$.

  \item[\termblanklong{parametric curve}]
    The set of ordered pairs $(x(t),\, y(t))$ traced in the $xy$-plane
    as $t$ \underline{\hspace{2.5cm}} through its domain.
    The curve has a \underline{\hspace{2.5cm}} (direction) determined
    by increasing $t$.

  \item[\termblanklong{start / end point}]
    When the domain is $a \le t \le b$, the curve begins at
    \underline{\hspace{2.5cm}} and ends at \underline{\hspace{2.5cm}}.
\end{description}
\end{vocabbox}

\vspace{0.08in}

%----------------------------------------------------------------------
% Hook
%----------------------------------------------------------------------
\begin{hookbox}
\textbf{Entry Question:} A drone's position at time $t$ seconds is
$x(t) = t^2 - 4$ (meters east of launch pad) and $y(t) = 2t$ (meters
north) for $0 \le t \le 4$.

\vspace{4pt}
\textbf{Q1.} What are the drone's coordinates when $t = 0$? When $t = 2$?
\writeline

\textbf{Q2.} Could you describe the drone's path with a single equation
$y = f(x)$? What information would that equation \emph{not} tell you?
\writelines{2}
\end{hookbox}

\vspace{0.08in}

%----------------------------------------------------------------------
% LO 4.1.A — Table and Graph
%----------------------------------------------------------------------
\begin{notesbox}{LO 4.1.A --- Parametric Functions: Table and Graph}

\textbf{Definition:} A parametric function $f(t) = (x(t),\, y(t))$
consists of \underline{\hspace{3cm}} component equations. Each value of
the parameter $t$ produces one ordered pair $(x(t),\, y(t))$ in the
$xy$-plane.

\vspace{8pt}
\textbf{Example:} Let $f(t) = (t + 1,\; t^2 - 3)$ for
$-2 \le t \le 2$.
Here $x(t) = t + 1$ and $y(t) = t^2 - 3$.

\vspace{6pt}
\textbf{Step 1 --- Build the table.} Evaluate each component at every
$t$-value in the domain.

\vspace{4pt}
\begin{center}
\renewcommand{\arraystretch}{1.6}
\begin{tabular}{|c|c|c|c|}
\hline
$t$ & $x(t) = t + 1$ & $y(t) = t^2 - 3$ & $(x,\; y)$ \\
\hline
$-2$ & \blank{1.8cm} & \blank{1.8cm} & \blank{2.4cm} \\
\hline
$-1$ & \blank{1.8cm} & \blank{1.8cm} & \blank{2.4cm} \\
\hline
$0$  & \blank{1.8cm} & \blank{1.8cm} & \blank{2.4cm} \\
\hline
$1$  & \blank{1.8cm} & \blank{1.8cm} & \blank{2.4cm} \\
\hline
$2$  & \blank{1.8cm} & \blank{1.8cm} & \blank{2.4cm} \\
\hline
\end{tabular}
\end{center}

\vspace{8pt}
\textbf{Step 2 --- Plot and connect.} Plot each $(x, y)$ ordered pair.
Then connect the points \emph{in order of increasing $t$} and draw
arrows to show the direction of the curve.

\vspace{6pt}
\begin{center}
\begin{tikzpicture}[scale=0.78]
  \draw[linegray, step=1cm] (-2,-4) grid (4,2);
  \draw[->] (-2.3,0) -- (4.5,0) node[right] {\small $x$};
  \draw[->] (0,-4.3) -- (0,2.5) node[above] {\small $y$};
  \foreach \x in {-2,-1,1,2,3,4} {
    \draw (\x, 0.06) -- (\x, -0.06);
    \node[below, font=\scriptsize] at (\x, 0) {$\x$};
  }
  \foreach \y in {-4,-3,-2,-1,1,2} {
    \draw (0.06, \y) -- (-0.06, \y);
    \node[left, font=\scriptsize] at (0, \y) {$\y$};
  }
\end{tikzpicture}
\end{center}

\vspace{4pt}
The start point is \underline{\hspace{2.5cm}} (at $t = -2$) and the
end point is \underline{\hspace{2.5cm}} (at $t = 2$).

\end{notesbox}

\vspace{0.08in}

%----------------------------------------------------------------------
% LO 4.1.A — Solving for t
%----------------------------------------------------------------------
\begin{notesbox}{LO 4.1.A --- Solving for the Parameter (Skill 1.A)}

\textbf{Key idea:} To find \emph{when} the curve reaches a specific
$x$- or $y$-value, set the component equation equal to that value and
solve for $t$.

\vspace{8pt}
\textbf{Example (using $f(t) = (t+1,\; t^2-3)$, $-2 \le t \le 2$):}

\vspace{4pt}
\textbf{Q:} At what point does the curve cross the $y$-axis
(where $x = 0$)?

Set $x(t) = 0$: \quad $t + 1 = 0 \;\Rightarrow\; t = $ \blank{1.2cm}

Find $y$: \quad $y\bigl($\blank{1cm}$\bigr) = $ \blank{1cm}${}^2 - 3 = $ \blank{1.5cm}

The curve crosses the $y$-axis at the point \blank{2.5cm}.

\vspace{10pt}
\textbf{Guided Practice.}\quad Use $f(t) = (t+1,\; t^2-3)$,
$-2 \le t \le 2$.

\vspace{4pt}
(a)\quad Find all values of $t$ in the domain where $y(t) = -2$.
Show your algebra.
\writelines{3}

(b)\quad What are the coordinates of all points on the curve where
$y = -2$?
\writeline

\end{notesbox}

\vspace{0.08in}

%----------------------------------------------------------------------
% Practice
%----------------------------------------------------------------------
\begin{practicebox}

\textbf{Practice.}\quad Let $g(t) = (2t - 3,\; 4 - t^2)$ for
$-1 \le t \le 2$.

\vspace{4pt}
(a)\quad Complete the table of values.

\vspace{4pt}
\begin{center}
\renewcommand{\arraystretch}{1.6}
\begin{tabular}{|c|c|c|c|}
\hline
$t$ & $x(t) = 2t - 3$ & $y(t) = 4 - t^2$ & $(x,\; y)$ \\
\hline
$-1$ & \blank{1.8cm} & \blank{1.8cm} & \blank{2.4cm} \\
\hline
$0$  & \blank{1.8cm} & \blank{1.8cm} & \blank{2.4cm} \\
\hline
$1$  & \blank{1.8cm} & \blank{1.8cm} & \blank{2.4cm} \\
\hline
$2$  & \blank{1.8cm} & \blank{1.8cm} & \blank{2.4cm} \\
\hline
\end{tabular}
\end{center}

\vspace{6pt}
(b)\quad What are the start and end points of the curve?
\writeline

(c)\quad Solve for $t$ when $x(t) = 1$. State the coordinates of
that point.
\writelines{2}

\end{practicebox}

\end{document}
